\newlabel{alg:exp_short}{{1}{25}{\textsc {Proof-Carrying Agentic Verification}}{algorithm.1}{}}
\newlabel{appendix:extended_proofs:cor_alpha}{{C.2}{30}{Certified operating-point corollary}{subsection.C.2}{}}
\newlabel{appendix:validation:controller}{{E.30}{53}{Controller calibration and broader sensitivity}{subsection.E.30}{}}
\newlabel{ass:dep_control}{{A2}{27}{Calibrated replay and control}{assumption.2}{}}
\newlabel{ass:integrity_replay}{{A1}{25}{Declared audit semantics and taxonomy coverage}{assumption.1}{}}
\newlabel{def:cert}{{D2}{3}{Replayable acceptance certificate}{definition.2}{}}
\newlabel{def:eg}{{D1}{3}{Agentic runtime graph}{definition.1}{}}
\newlabel{def:grounded}{{D3}{3}{Checkable grounded acceptance}{definition.3}{}}
\newlabel{def:indep}{{D4}{3}{$(\delta ,\kappa ,\zeta )$-separated support paths}{definition.4}{}}
\newlabel{eq:resp}{{1}{4}{Support and interventions}{equation.1}{}}
\newlabel{fig:app-audit-sampling}{{15}{35}{\footnotesize \Derived . \textbf {Audit-design scorecard.} \textbf {A}: worst-stratum coverage (higher is better). \textbf {B}: certified envelope (lower is tighter). \textbf {C}: weighting-adjusted effective sample size. \textbf {D}: violation and false-certification rates (lower is better). All panels use the same declared sampling design; effective sample size is not itself a safety guarantee}{figure.caption.46}{}}
\newlabel{fig:app-injection}{{17}{36}{\footnotesize \Direct /\Derived \ . \textbf {Injection and common-mode stress.} \textbf {A}: attack success, accepted attack, and detection under the declared challenge; lower accepted attack and higher detection are preferred. \textbf {B}: supplied dependence-inflation bounds for matched injection regimes; larger inflation weakens redundancy claims and is not a correlation coefficient}{figure.caption.50}{}}
\newlabel{fig:app-responsibility}{{18}{37}{\footnotesize \Direct \ . \textbf {Replay-interventional diagnosis under closed, mixed, and open-set failures.} \textbf {A}: supplied top-1/top-3 diagnostic accuracy. \textbf {B}: unresolved cases remain explicit rather than being forced into a component label. \textbf {C}: median responsibility margin measures diagnostic separation. These are replay-relative diagnostics, not real-world causal-identification claims}{figure.caption.53}{}}
\newlabel{fig:app-theory-geometry}{{4}{7}{\footnotesize \Protocol /\Derived . \textbf {Capability scope and theoretical envelopes.} \textbf {A}: FULL, PARTIAL, and N/A denote declared capability scope; recomputability requires an externally replayable decision contract. \textbf {B}: contours evaluate the declared dependence/redundancy envelope; the marked point is the protocol setting and the closed region fails the gate. \textbf {C}: audit probes reduce sampling uncertainty, while uncovered deployment mass enlarges the bound. The shading gives empirical 95\% Wilson confidence bands around the supplied residual rate}{figure.caption.5}{}}
\newlabel{fig:app-witnesses}{{16}{36}{\footnotesize \Direct \ . \textbf {Separating witnesses.} Each challenge targets one PCG-MAS conjunct while the remaining channels are intended to pass. Cells show acceptance rates, so low acceptance in the targeted PCG-MAS row is desirable. FULL/PARTIAL denote declared comparator scope; hatched N/A is out of scope, not rejection}{figure.caption.48}{}}
\newlabel{fig:app_replay_modes}{{8}{18}{\footnotesize \Protocol . \textbf {Snapshot replay versus fresh-environment drift.} Snapshot replay reconstructs committed evidence, hashes, and tool outputs to test reproducibility; fresh execution reissues external interactions and is monitored as $\mathsf {DriftFail}$. The two answer different questions and are not merged into one acceptance bit}{figure.caption.16}{}}
\newlabel{fig:app_support_paths}{{9}{20}{\footnotesize \Protocol . \textbf {Separated support paths and delegation-aware diagnosis.} Candidate branches expose provenance roots, evidence, tools/parsers, and delegates. The gate checks declared overlap/separation constraints before redundancy is credited; replay-interventional responsibility is then computed only over the committed graph}{figure.caption.18}{}}
\newlabel{fig:r1_to_r4_combined}{{5}{8}{\footnotesize \Derived . \textbf {Mechanism and falsification atlas.} \textbf {A}: accepted-risk change under component ablation; larger positive change means the removed component mattered more. \textbf {B}: supplied risk--coverage curves and the seven observed aggregate operating points; lower risk and higher coverage are preferred, and the selected PCG-MAS point is distinguished from comparator points. \textbf {C}: the supplied dataset-level selectivity and verification components satisfy the exact identity $\Delta =S+V$ and match Table~\ref {tab:table_25_protocol_sv}. \textbf {D}: accepted-risk reduction versus the best applicable comparator; positive favors PCG-MAS}{figure.caption.7}{}}
\newlabel{fig:robustness_atlas}{{6}{9}{\footnotesize \Direct /\Derived . \textbf {Adversarial and deployment-shift robustness.} \textbf {A}: attack success versus accepted attack across the declared injection regimes. \textbf {B}: detection together with the dependence-inflation bound $\widehat {\rho }_{\mathrm {UCB}}$, exposing common-mode regimes where redundancy credit weakens. \textbf {C}: bound-violation rate versus shift severity under static calibration, monitored gating, and rolling recalibration. \textbf {D}: utility, audit coverage, and alarm behavior across the same shift range. Lower accepted attack and violation rates are preferred; higher detection and maintained utility/coverage are preferred}{figure.caption.8}{}}
\newlabel{lem:sv_identity}{{C.2}{29}{Exact selectivity--verification decomposition}{lemma.C.2}{}}
\newlabel{lem:trace_separation}{{C.1}{26}{Trace decomposition}{lemma.C.1}{}}
\newlabel{thm:audit_converse}{{C.2}{27}{Confidence-rate converse and unavoidable uncovered mass}{theorem.C.2}{}}
\newlabel{thm:audit_decomp}{{C.1}{27}{Finite-sample contract-audit envelope}{proposition.C.1}{}}
\newlabel{thm:commonmode}{{C.4}{28}{Common-mode floor}{theorem.C.4}{}}
\newlabel{thm:irredundancy}{{C.1}{25}{Irredundancy of the four certificate conjuncts}{theorem.C.1}{}}
\newlabel{thm:restricted_tv}{{C.5}{28}{Restricted-TV transfer and domain-classifier identity}{theorem.C.5}{}}
\newlabel{thm:rho_renyi}{{C.3}{28}{All-fail likelihood ratio and R\'enyi envelope}{theorem.C.3}{}}
\newlabel{thm:threshold_refusal}{{C.2}{29}{Responsibility ranking and controller consequence}{proposition.C.2}{}}
